\chapter{Externalities and Public Goods: The Formal Theory}\label{ch:b2:externalities-and-public-goods-theory}

Burning coal with sulphur in it releases sulphur dioxide, which forms fine particles that people breathe and acid that
falls with the rain. The plant that burns the coal does not pay for that harm. From 1995 the most polluting American power plants needed a
permit, an allowance, for every ton they emitted, and from 2000 all large plants shared a national cap. The chart
shows what happened to emissions.

\begin{openingchart}{Sulphur-dioxide emissions in the United States, million tonnes a year, 1970--2023: electricity and heat
  production, and all other sources. Dashed lines: the start of the allowance programme's first phase (1995) and of
  its national cap (2000).}{Community Emissions Data System (CEDS), emissions by country and sector}
\begin{tikzpicture}
\begin{axis}[omrecord, width=0.82\linewidth, height=0.42\linewidth, xmin=1970, xmax=2023, ymin=0, ymax=19,
    xlabel={Year}, ylabel={Million tonnes}, xtick={1970,1980,1990,2000,2010,2020},
    xticklabel style={/pgf/number format/1000 sep=}, legend style={at={(0.5,-0.3)}, anchor=north,
    legend columns=2, /tikz/every even column/.append style={column sep=10pt}}]
  \addplot[omDef, very thick] table[x=year, y=power]
    {figdata/bachelor-2/18-externalities-and-public-goods-theory/out/so2.csv};
  \addplot[omThm, thick, dashed] table[x=year, y=other]
    {figdata/bachelor-2/18-externalities-and-public-goods-theory/out/so2.csv};
  \draw[black!50, dashed] (axis cs:1995,0) -- (axis cs:1995,19);
  \draw[black!50, dashed] (axis cs:2000,0) -- (axis cs:2000,19);
  \legend{Electricity and heat, All other sources}
\end{axis}
\end{tikzpicture}
\end{openingchart}

% figdata: so2_summary (power_1990 14.76, power_1995 11.91, power_2000 11.17, power_2010 4.88, power_last 0.79, power_fall_pct 95, last 2023)
In this series, emissions from electricity and heat fell from 14.8 million tonnes in 1990 to 11.9 in 1995, 11.2 in 2000,
4.9 in 2010 and 0.8 in 2023, a fall of 95 per cent. Emissions from all other sources had fallen earlier. Later rules on the
same pollutant, which required further cuts state by state and limited trading between states, also lowered emissions
(Schmalensee and Stavins 2013): the chart alone cannot say how much of the fall each rule caused.

This chapter states the theory behind such programmes: why an external effect makes a \omterm{def:b2:exchange-economies:competitive-equilibrium}{competitive equilibrium}
inefficient, the three classical remedies (a tax, permits and bargaining), how to choose between a price and a quantity
when costs are uncertain, and the parallel theory of \omterm{def:b2:externalities-and-public-goods-theory:public-good}{public goods}, including a mechanism that makes people state their
values truthfully. It ends with the ways public remedies themselves fail.

\begin{recall}
From the school volumes: externalities, \omterm{def:b2:externalities-and-public-goods-theory:public-good}{public goods}, the \omterm{def:b2:externalities-and-public-goods-theory:corrective-tax}{corrective tax} and the tragedy of the commons. From
\cref{ch:b2:exchange-economies,ch:b2:welfare-theorems}: \omterm{def:b2:exchange-economies:competitive-equilibrium}{competitive equilibrium} and the first welfare theorem. From
\cref{ch:b2:static-games}: voluntary contributions to a \omterm{def:b2:externalities-and-public-goods-theory:public-good}{public good}.
\end{recall}

\begin{notation}[Symbols of this chapter]\label{not:b2:externalities-and-public-goods-theory:symbols}
$E$ is a quantity of emissions; $\mathit{MB}(E)$ the polluter's marginal benefit of emitting, equal to its marginal cost of
abatement; $\mathit{MD}(E)$ the marginal damage to others. $G$ is a quantity of a \omterm{def:b2:externalities-and-public-goods-theory:public-good}{public good}, $v_i(G)$ person $i$'s benefit from it, and
$c(G)$ its cost.
\end{notation}

\section{An external effect in a competitive economy}\label{sec:b2:externalities-and-public-goods-theory:inefficiency}

\begin{definition}[Externality]\label{def:b2:externalities-and-public-goods-theory:externality}
An \emph{externality}\index{externality} is an effect of one agent's consumption or production on another agent's
utility or production that does not pass through prices.
\end{definition}

The first welfare theorem (\cref{ch:b2:welfare-theorems}) assumed that each agent's utility or output depends only on
what that agent buys and sells. An \omterm{def:b2:externalities-and-public-goods-theory:externality}{externality} breaks the assumption, and with it the theorem. Take a mill whose
emissions $E$ are a by-product of its output, so that it gains $\mathit{MB}(E)$, decreasing, from a further unit of emissions; and a
fishery downstream that loses $\mathit{MD}(E)$, increasing, from it. Utility is quasi-linear in money, so surplus can be added.

\begin{proposition}[Competitive equilibrium with an externality is inefficient]\label{prop:b2:externalities-and-public-goods-theory:inefficient}
If the mill chooses emissions to maximise its own profit and $\mathit{MD}(E) > 0$, it emits more than the level $E^*$ that maximises total
surplus, defined by $\mathit{MB}(E^*) = \mathit{MD}(E^*)$.
\end{proposition}

\begin{proof}
The mill emits until its own marginal benefit is zero: $\mathit{MB}(E_0) = 0$. Total surplus $\int_0^E (\mathit{MB} - \mathit{MD})$ is maximised where $\mathit{MB} = \mathit{MD}$. At $E_0$,
$\mathit{MB} - \mathit{MD} = -\mathit{MD}(E_0) < 0$, so a small cut in emissions raises surplus. Since $\mathit{MB} - \mathit{MD}$ is decreasing, it is zero at a unique $E^*$ below
$E_0$.
\end{proof}

Take $\mathit{MB}(E) = 120 - 2E$ and $\mathit{MD}(E) = E$, in thousands of dalens per tonne. The mill emits 60 tonnes; the efficient level is 40, where $120 -
2E = E$. Between 40 and 60 each tonne does more harm than good, and the surplus lost is the triangle between the two curves, $\int_{40}^{60}(3E - 120)\,dE = 600$.

\section{A tax and permits}\label{sec:b2:externalities-and-public-goods-theory:tax}

\begin{definition}[Corrective tax]\label{def:b2:externalities-and-public-goods-theory:corrective-tax}
A \emph{corrective tax}\index{corrective tax} is a tax on an activity equal to the marginal external damage it causes at the
efficient level, so that the agent's private cost equals the social cost.
\end{definition}

\begin{definition}[Tradable emission permit]\label{def:b2:externalities-and-public-goods-theory:permit}
A \emph{tradable emission permit}\index{tradable emission permit} is a right to emit one unit of a pollutant in a period,
which its holder may use, sell or keep; the government fixes the total number.
\end{definition}

A tax of $t = \mathit{MD}(E^*) = 40$ a tonne makes the mill emit where $\mathit{MB}(E) = 40$, which is 40 tonnes: the efficient level. Alternatively the
government can issue 40 permits. If the mill must buy them, their price settles where the mill's marginal benefit of
the fortieth tonne equals the price, again 40. With one polluter the two instruments are the same, one fixing the price and
letting the quantity follow, the other the reverse. With many polluters, both have the property that matters most: every
polluter abates until its marginal cost of abatement equals the same price, so a given total is reached at least cost
(\cref{exo:b2:externalities-and-public-goods-theory:6}). Who pays differs: a tax and auctioned permits raise revenue; permits
given away free hand their value to the polluters, without changing the outcome.

\begin{omfigure}
\begin{tikzpicture}[econ]
\begin{axis}[omrecord, grid=none, xmin=0, xmax=66, ymin=0, ymax=130, width=0.6\linewidth,
    height=0.45\linewidth, xlabel={$E$: tonnes emitted}, ylabel={Thousands of dalens a tonne},
    xtick={0,40,60}, ytick={40,120}, legend style={at={(1.03,0.5)}, anchor=west}]
  \addplot[fill=omDef!15, draw=none, forget plot] coordinates {(40,40) (60,60) (60,0)} -- cycle;
  \addplot[omDef, very thick, domain=0:60, samples=2] {120 - 2*x};
  \addplot[omThm, very thick, dashed, domain=0:66, samples=2] {x};
  \addplot[black!60, thin, domain=0:66, samples=2] {40};
  \node[eqpoint] at (axis cs:40,40) {};
  \node[font=\footnotesize] at (axis cs:55,22) {loss};
  \legend{Mill's marginal benefit, Marginal damage, Tax of 40}
\end{axis}
\end{tikzpicture}

{\small A polluter and a \omterm{def:b2:externalities-and-public-goods-theory:corrective-tax}{corrective tax}. Unregulated, the mill emits until its marginal benefit is zero, at 60 tonnes; the
efficient level is 40, where marginal benefit equals marginal damage. A tax of 40 a tonne makes the mill stop there. The
shaded triangle is the surplus lost without the tax.}
\end{omfigure}

\begin{record}[The sulphur-dioxide permit market]
% ledger: b2:ss2013-design, b2:ss2013-performance, b2:ss2013-benefits
The 1990 amendments to the United States Clean Air Act aimed to cut annual sulphur-dioxide emissions by 10 million tons from
their 1980 level. From 1995 the 263 most polluting coal-fired units needed allowances; from 2000 about 3{,}200 units shared a
cap of 8.95 million tons a year, half the 1980 level. Allowances were given free, and could be traded or banked. Emissions
from power plants fell 36 per cent between 1990 and 2004 while coal-fired generation rose 25 per cent, and the cap was reached
in 2007. With continuous monitoring and a fine of 2{,}000 dollars a ton, compliance was nearly complete. Estimates put the
saving against rules that prescribed the means of abatement at between 15 and 90 per cent, and the annual benefits, mostly
from fewer deaths from fine particles, at 59 to 116 billion dollars against costs of 0.5 to 2 billion (Schmalensee and Stavins
2013).
\end{record}

\begin{measure}[How emissions are counted]
% ledger: b2:hoesly2018, b2:ss2013-performance
Most emissions are not measured at the chimney but estimated: the fuel burned in each sector is multiplied by an emission
factor, the tonnes of pollutant per unit of fuel given the technology, and the results are reconciled with national
inventories. The series in the opening chart is built this way, consistently for all years and countries; its recent
years are the most uncertain (Hoesly et al.\ 2018). Large American power plants under the permit programme measured their
emissions continuously, which made a cap enforceable (Schmalensee and Stavins 2013). Different series of the same
emissions differ in what they cover and how they estimate, so levels should be compared only within one series.
\end{measure}

\begin{model}{Corrective tax and permits}
\assumes A known marginal damage; polluters that minimise cost; emissions that can be measured; a government able to set
and keep a tax rate or a cap.
\predicts Emissions fall to the level where marginal abatement cost equals the tax or permit price, at least total cost
for that level; abatement moves to the polluters who can abate most cheaply.
\tested The United States sulphur-dioxide market: the cap was met early and at well below the cost of prescriptive rules
(Schmalensee and Stavins 2013).
\breaks When damage depends on where the pollutant is emitted, so that a tonne is not a tonne everywhere; and when the
government changes the rules, as in the collapse of the same market after 2008.
\end{model}

\section{Bargaining: the Coase theorem}\label{sec:b2:externalities-and-public-goods-theory:coase}

\begin{definition}[Transaction cost]\label{def:b2:externalities-and-public-goods-theory:transaction-cost}
A \emph{transaction cost}\index{transaction cost} is a cost of making an exchange other than the price: finding the
other party, negotiating, writing and enforcing the agreement.
\end{definition}

\begin{definition}[Coase theorem]\label{def:b2:externalities-and-public-goods-theory:coase}
The \emph{Coase theorem}\index{Coase theorem} states that if property rights over an external effect are clearly
assigned, \omterm{def:b2:externalities-and-public-goods-theory:transaction-cost}{transaction costs} are zero and utility is quasi-linear, bargaining between the parties reaches the efficient
level of the effect whichever party holds the right; the assignment affects only who pays whom.
\end{definition}

\begin{proposition}[The Coase theorem]\label{prop:b2:externalities-and-public-goods-theory:coase}
With quasi-linear utility and costless bargaining that exhausts the gains from trade, the mill and the fishery agree on $E^*$
whether the mill has the right to pollute or the fishery the right to clean water.
\end{proposition}

\begin{proof}
With quasi-linear utility, money moves surplus one for one between the parties, so any agreement that does not maximise
their joint surplus $\int_0^E (\mathit{MB} - \mathit{MD})$ can be improved for both by changing $E$ toward the maximiser and dividing the gain. Bargaining
that exhausts the gains therefore ends at $E^*$, which by \cref{prop:b2:externalities-and-public-goods-theory:inefficient} does not depend on who
holds the right. The right fixes the starting point, $E_0$ or 0, and so the direction of the payment: the fishery pays the
mill to cut from $E_0$, or the mill pays the fishery for permission to emit up to $E^*$.
\end{proof}

In the example, if the mill has the right, cutting from 60 to 40 tonnes costs it $\int_{40}^{60}(120 - 2E)\,dE = 400$ and saves the fishery
$\int_{40}^{60}E\,dE = 1{,}000$; any payment between 400 and 1{,}000 makes both better off. The theorem's lesson is less that bargaining
solves externalities than that the obstacles to bargaining are the problem: many parties, private information about
costs and damages, and the cost of enforcing an agreement. Where these are small, private agreements appear.

\begin{record}[Bees and orchards]
% ledger: b2:cheung1973
Pollination by bees is a textbook example of an \omterm{def:b2:externalities-and-public-goods-theory:externality}{externality} between beekeepers and fruit growers. A survey in the state of
Washington in spring 1972, of nine beekeepers with about 10{,}000 colonies, found it handled by contract. Hives are moved by
truck from farm to farm, about 2.2 crops a year. Growers of crops that need pollination but yield little honey, such as
cherries, pay beekeepers a pollination fee; for crops that yield honey but need no pollination, such as mint, beekeepers pay
the farmer a rent, mostly in honey. Where the hives on one orchard also pollinate the neighbours', a custom replaces
contracts among the many neighbours: each grower keeps or hires as many hives per acre as the orchards around him (Cheung
1973).
\end{record}

\begin{model}{Coase bargaining}
\assumes Clearly assigned rights; few parties who know one another's costs and damages; agreements that can be enforced;
quasi-linear utility.
\predicts The efficient level of the external effect whatever the assignment of rights; payments in the direction the
rights determine.
\tested Pollination contracts between beekeepers and growers, with fees and rents that run in the direction the crop
implies (Cheung 1973).
\breaks With many parties, such as the residents downwind of a power plant: no one can gather them to bargain, and each
would hope the others pay.
\end{model}

\section{Prices or quantities?}\label{sec:b2:externalities-and-public-goods-theory:weitzman}

When the regulator knows both curves, a tax and a cap are equivalent. Usually it does not know the polluter's cost of
abatement, which the polluter knows better. Let $\mathit{MB}(E) = a + \theta - bE$, with $\theta$ a shock to abatement cost of mean zero that the
polluter sees and the regulator does not, and $\mathit{MD}(E) = dE$. The regulator must set a tax $t$ or a quota $\bar E$ before $\theta$ is known. It
sets both for the average case: $\bar E = a/(b + d)$ and $t = d\bar E$.

\begin{proposition}[The slopes rule]\label{prop:b2:externalities-and-public-goods-theory:slopes}
The expected welfare loss of the tax is $d^2/b^2$ times that of the quota. A tax is better when marginal damage is flatter
than marginal abatement cost ($d < b$), a quota when it is steeper.
\end{proposition}

\begin{proof}
The efficient level is $E^*(\theta) = (a + \theta)/(b + d)$. The loss from any $E$ is the triangle between the curves, $\tfrac12(b + d)(E - E^*)^2$. Under the
quota, $E - E^* = -\theta/(b + d)$, so the loss is $\theta^2/(2(b + d))$. Under the tax, the polluter emits where $a + \theta - bE = t$, so $E = \bar E + \theta/b$ and $E - E^* = \theta/b -
\theta/(b + d) = \theta d/(b(b + d))$; the loss is $\theta^2 d^2/(2b^2(b + d))$. Taking expectations, the ratio of the two is $d^2/b^2$.
\end{proof}

A tax fixes the marginal cost of abatement and lets emissions vary with the shock; a quota fixes emissions and lets the
marginal cost vary. If damage barely changes with a few more tonnes, varying emissions costs little, and a tax is better;
if damage rises steeply, as near a threshold, fixing the quantity is better.

\begin{omfigure}
\begin{tikzpicture}[econ]
\begin{axis}[omrecord, grid=none, xmin=8, xmax=31, ymin=0, ymax=60, width=0.6\linewidth,
    height=0.45\linewidth, xlabel={$E$: tonnes emitted}, ylabel={Thousands of dalens a tonne},
    xtick={15,20,25}, ytick={30}, legend style={at={(1.03,0.5)}, anchor=west}]
  \addplot[omDef, very thick, domain=10:30, samples=2] {105 - 3*x};
  \addplot[omDef, thick, dashed, domain=8:25, samples=2] {75 - 3*x};
  \addplot[omThm, very thick, domain=8:31, samples=2] {1.5*x};
  \addplot[black!60, thin, domain=8:31, samples=2] {30};
  \draw[black!60, thin] (axis cs:20,0) -- (axis cs:20,60);
  \node[eqpoint] at (axis cs:23.333,35) {};
  \node[eqpoint] at (axis cs:16.667,25) {};
  \node[anchor=south west, font=\footnotesize] at (axis cs:20,55) {quota};
  \node[anchor=south, font=\footnotesize] at (axis cs:10,30.5) {tax};
  \legend{Marginal benefit (high cost), Marginal benefit (low cost), Marginal damage}
\end{axis}
\end{tikzpicture}

{\small Prices against quantities, with the numbers of the weekend problem. Abatement cost may be high or low. The efficient
emissions (dots) are 23.3 and 16.7 tonnes. The quota holds emissions at 20 in both cases; the tax of 30 lets them move to 25
and 15, closer to the efficient levels because marginal damage is flatter than marginal benefit.}
\end{omfigure}

The sulphur-dioxide market was a quantity instrument, and its price did the varying. It rose from 273 dollars a ton at the
2004 auction to 703 at the 2005 auction, and passed 1{,}200 dollars in the market later that year, on expectations of a
tighter cap. It fell to 70 at the 2009 auction and 56 cents at the 2012 one, as
the rules changed. The chart shows the auction prices that Schmalensee and Stavins quote.

\begin{omchart}{Spot price of a sulphur-dioxide allowance at the Environmental Protection Agency's annual auction, dollars a
  ton, logarithmic scale, the five years quoted by the source.}{Schmalensee and Stavins (2013)}
\begin{tikzpicture}
\begin{axis}[omrecord, ymode=log, width=0.6\linewidth, height=0.38\linewidth, xmin=2003, xmax=2013, ymin=0.3,
    ymax=2000, xtick={2004,2006,2008,2010,2012}, xticklabel style={/pgf/number format/1000 sep=},
    ytick={1,10,100,1000}, yticklabels={1,10,100,1000}, xlabel={Year of auction}, ylabel={Dollars a ton},
    nodes near coords, point meta=explicit symbolic, nodes near coords style={font=\footnotesize, anchor=west}]
  \addplot[omDef, thick, mark=*, mark size=2pt] table[x=year, y=price, meta=price]
    {figdata/bachelor-2/18-externalities-and-public-goods-theory/out/allowance.csv};
\end{axis}
\end{tikzpicture}
\end{omchart}

\begin{model}{Prices against quantities}
\assumes A regulator uncertain about abatement costs, who must fix a tax or a quota in advance; linear marginal curves.
\predicts A tax gives certain prices and uncertain emissions, a quota the reverse; the tax has the lower expected loss when
marginal damage is flatter than marginal abatement cost.
\tested The slopes rule has no direct test reported here: a test would compare realised losses under a tax and under a
cap for the same pollutant and the same cost shocks. The sulphur-dioxide cap illustrates the price risk of a quantity
instrument: emissions stayed within the cap while the allowance price ranged from under a dollar to over 1{,}200 dollars
a ton in the market (Schmalensee and Stavins 2013).
\breaks When the uncertainty is about the rules themselves rather than about costs: the same price swings came mostly from
proposed and annulled regulations.
\end{model}

\section{Public goods and the Samuelson condition}\label{sec:b2:externalities-and-public-goods-theory:samuelson}

\begin{definition}[Public good]\label{def:b2:externalities-and-public-goods-theory:public-good}
A \emph{public good}\index{public good} is a good that is non-rival, so that one person's use does not reduce what
others can use, and non-excludable, so that no one can be kept from using it.
\end{definition}

\begin{definition}[Club good]\label{def:b2:externalities-and-public-goods-theory:club-good}
A \emph{club good}\index{club good} is a good that is non-rival, at least until it is crowded, but excludable, so that
its use can be charged for.
\end{definition}

\begin{definition}[Samuelson condition]\label{def:b2:externalities-and-public-goods-theory:samuelson}
The \emph{Samuelson condition}\index{Samuelson condition} states that the efficient quantity of a \omterm{def:b2:externalities-and-public-goods-theory:public-good}{public good} equates
the sum of all users' marginal benefits to its marginal cost.
\end{definition}

\begin{proposition}[Efficient provision of a public good]\label{prop:b2:externalities-and-public-goods-theory:samuelson}
With quasi-linear utilities $v_i(G) + x_i$, $v_i$ increasing and concave, and a convex cost $c(G)$ paid from the $x_i$, a quantity $G^* > 0$ is
efficient if and only if $\sum_i v_i'(G^*) = c'(G^*)$.
\end{proposition}

\begin{proof}
With quasi-linear utility, an \omterm{def:b2:exchange-economies:allocation}{allocation} is efficient if and only if it maximises total surplus $\sum_i v_i(G) - c(G)$: if it did not,
a change in $G$ with suitable transfers would make everyone better off. The surplus is concave, so it is maximised where
its derivative is zero, $\sum_i v_i'(G) = c'(G)$.
\end{proof}

For a private good, efficiency requires each person's marginal benefit to equal marginal cost; for a \omterm{def:b2:externalities-and-public-goods-theory:public-good}{public good}, the
sum, because one unit serves everyone. Take three households on a street that could light it with lamps: their marginal
benefits from the $G$-th lamp are $12 - G$, $10 - G$ and $8 - G$ dalens a month, and a lamp costs 9. The sum is $30 - 3G$, so the efficient
number is 7.

\begin{omfigure}
\begin{tikzpicture}[econ]
\begin{axis}[omrecord, grid=none, xmin=0, xmax=10.5, ymin=0, ymax=31, width=0.6\linewidth,
    height=0.45\linewidth, xlabel={$G$: lamps}, ylabel={Dalens a month},
    xtick={3,7,10}, ytick={9,12,30}, legend style={at={(1.03,0.5)}, anchor=west}]
  \addplot[omDef, thin, domain=0:10, samples=2] {12 - x};
  \addplot[omDef, thin, dashed, domain=0:10, samples=2] {10 - x};
  \addplot[omDef, thin, dotted, domain=0:8, samples=2] {8 - x};
  \addplot[omThm, very thick, domain=0:8, samples=2] {30 - 3*x};
  \addplot[black!60, thin, domain=0:10.5, samples=2] {9};
  \node[eqpoint] at (axis cs:7,9) {};
  \node[eqpoint] at (axis cs:3,9) {};
  \node[anchor=south west, font=\footnotesize] at (axis cs:7,9.3) {efficient};
  \draw[guide] (axis cs:3,0) -- (axis cs:3,9);
  \node[anchor=north west, font=\footnotesize, fill=white, inner sep=1pt] at (axis cs:3.1,3) {voluntary};
  \legend{Household 1, Household 2, Household 3, Sum, Cost of a lamp}
\end{axis}
\end{tikzpicture}

{\small Vertical summation. Each lamp lights the whole street, so the marginal benefits of the three households are added
at each quantity. The sum meets the cost of a lamp at 7 lamps; left to themselves, the households buy 3.}
\end{omfigure}

\section{Voluntary provision falls short}\label{sec:b2:externalities-and-public-goods-theory:voluntary}

If each household buys lamps for itself, taking the others' purchases as given (\cref{ch:b2:static-games}), household 1 buys
lamps until its own marginal benefit equals the cost: $12 - G = 9$, so $G = 3$. At 3 lamps the others' marginal benefits, 7 and 5, are
below 9, so they buy none. The street gets 3 lamps instead of 7, and the surplus lost is $\int_3^7 (21 - 3G)\,dG = 24$ dalens a month. Each
household counts only its own benefit, ignoring what its lamps give the others: provision of a \omterm{def:b2:externalities-and-public-goods-theory:public-good}{public good} is an
\omterm{def:b2:externalities-and-public-goods-theory:externality}{externality} in reverse.

\section{Personalised prices}\label{sec:b2:externalities-and-public-goods-theory:lindahl}

\begin{definition}[Personalised prices]\label{def:b2:externalities-and-public-goods-theory:lindahl}
\emph{Personalised (Lindahl) prices}\index{personalised prices} are prices per unit of a \omterm{def:b2:externalities-and-public-goods-theory:public-good}{public good}, one for each
person, that sum to its marginal cost and at which every person demands the same quantity.
\end{definition}

If household 1 paid 5 a lamp, household 2 paid 3 and household 3 paid 1, each would want exactly 7 lamps, since each one's marginal
benefit at 7 equals its price, and the prices sum to the cost, 9. \omterm{def:b2:externalities-and-public-goods-theory:lindahl}{Personalised prices} reach the efficient quantity as
ordinary prices do for private goods. But they need each person's marginal benefit, which each has every reason to
understate: a household that claims to value lamps less pays less and still enjoys the street lit by the others.

\section{A mechanism that makes people tell the truth}\label{sec:b2:externalities-and-public-goods-theory:pivot}

\begin{definition}[Pivot mechanism]\label{def:b2:externalities-and-public-goods-theory:pivot}
The \emph{pivot mechanism}\index{pivot mechanism} decides on a public project from reported values, carrying it out if
the reported net values sum to zero or more, and charges each person whose report changes the decision the net loss
that the change imposes on everyone else.
\end{definition}

\begin{proposition}[Truth-telling is a dominant strategy]\label{prop:b2:externalities-and-public-goods-theory:pivot}
In the \omterm{def:b2:externalities-and-public-goods-theory:pivot}{pivot mechanism}, reporting one's true net value is a \omterm{def:b2:static-games:best-response}{best response} to any reports by the others.
\end{proposition}

\begin{proof}
Let person $i$ have true net value $u_i$ and let $S$ be the sum of the others' reported net values. If the project is carried out,
$i$ gets $u_i$, and pays $-S$ if $S < 0$, since then she changed the decision; so her payoff is $u_i + \min(S, 0)$. If it is not, she gets 0, and pays $S$ if
$S \ge 0$: her payoff is $-\max(S, 0)$. Subtracting $\min(S, 0)$ from both, she compares $u_i + S$ with 0. She prefers the project exactly when $u_i + S \ge 0$. Reporting
$u_i$ makes the mechanism build exactly when $u_i + S \ge 0$, which is her preference, whatever $S$ is. No other report can do better, and
some do worse.
\end{proof}

A bridge costs 300 dalens, shared equally by three residents who value it at 250, 60 and 40. The listing
runs the mechanism and checks, by trying every report from 0 to 400, that no resident gains by lying.

\omcode{code/bachelor-2/18-externalities-and-public-goods-theory/listings/pivot.py}{1}{33}{The \omterm{def:b2:externalities-and-public-goods-theory:pivot}{pivot mechanism} for
a shared bridge, with a check of every alternative report.}\label{lst:b2:externalities-and-public-goods-theory:pivot}

\omoutput{code/bachelor-2/18-externalities-and-public-goods-theory/outputs/pivot.txt}{The bridge is built; only resident 1
is pivotal, and pays 100; no resident gains from any other report.}\label{out:b2:externalities-and-public-goods-theory:pivot}

% code: listings/pivot.py
Net values are 150, $-40$ and $-60$, summing to 50, so the bridge is built. Without resident 1 the others' net values sum to $-100$: she changes
the decision, and pays the 100 that the other two lose. The mechanism has a cost: the 100 cannot be handed back to the
residents, since a rebate that depended on the reports would change their incentives, so it must be wasted or paid to
someone outside. A mechanism with dominant-strategy truth-telling and an efficient decision cannot, in general, balance its
budget.

\section{Government failure beside market failure}\label{sec:b2:externalities-and-public-goods-theory:government}

Every remedy in this chapter is chosen and run by a government that has its own information problems and incentives. The
sulphur-dioxide cap was not set where marginal benefit met marginal damage; no credible estimates of the benefits existed,
and the level was chosen near where abatement costs began to rise steeply and to win political support. Its later
collapse came from the government: the law gave the regulator no power to tighten the cap as evidence on health damage
grew, a replacement rule was annulled by a court, and the rules that followed restricted trading between states, so that
demand for allowances all but vanished (Schmalensee and Stavins 2013). A regulator that overestimates damage, as in
\cref{exo:b2:externalities-and-public-goods-theory:11}, imposes losses as a missing regulator does. The theory of how
public choices are made, and how they fail, is the subject of the Year 3 volume.

\section{Exercises}\label{sec:b2:externalities-and-public-goods-theory:exercises}

\begin{exercise}[$\star$]\label{exo:b2:externalities-and-public-goods-theory:1}
Read the opening chart. By how much did emissions from electricity and heat fall from 1990 to 2023? Why does the chart alone
not show how much of the fall the permit programme caused?
\end{exercise}

\begin{exercise}[$\star$]\label{exo:b2:externalities-and-public-goods-theory:2}
Let $\mathit{MB}(E) = 100 - 2E$ and $\mathit{MD}(E) = 2E$. Find unregulated emissions, the efficient level, the \omterm{def:b2:externalities-and-public-goods-theory:corrective-tax}{corrective tax} and the surplus lost
without it.
\end{exercise}

\begin{exercise}[$\star$]\label{exo:b2:externalities-and-public-goods-theory:3}
Classify as public, club, common-pool or private: a coastal lighthouse; a subscription music service; a fishing ground open
to all; a sandwich. Give the property that decides each.
\end{exercise}

\begin{exercise}[$\star$]\label{exo:b2:externalities-and-public-goods-theory:4}
In the bridge example, change resident 2's value to 160. Is the bridge built, and who pays a tax? Then return to 60 and
raise the bridge's cost to 330: who pays, and how much?
\end{exercise}

\begin{exercise}[$\star\star$]\label{exo:b2:externalities-and-public-goods-theory:5}
In the example of \cref{sec:b2:externalities-and-public-goods-theory:inefficiency}, let the fishery hold the right to clean
water. What is the mill's gain from emitting 40 tonnes, the fishery's damage, and the range of payments that both accept?
\end{exercise}

\begin{exercise}[$\star\star$]\label{exo:b2:externalities-and-public-goods-theory:6}
Two firms each emit 60 tonnes unregulated; firm A's marginal abatement cost is $120 - 2E_A$, firm B's $60 - E_B$. A cap of 50 tonnes is
imposed. Compare the total abatement cost when each firm gets 25 permits and may not trade with the cost when they may
trade. Find the permit price.
\end{exercise}

\begin{exercise}[$\star\star$]\label{exo:b2:externalities-and-public-goods-theory:7}
Add a fourth household with marginal benefit $6 - G$ to the street, and let a lamp cost 12. Find the efficient number of lamps,
the \omterm{def:b2:externalities-and-public-goods-theory:lindahl}{personalised prices}, the number bought voluntarily and the surplus lost.
\end{exercise}

\begin{exercise}[$\star\star$]\label{exo:b2:externalities-and-public-goods-theory:8}
Predict, then compute: with $\mathit{MB}(E) = 90 + \theta - 3E$, $\theta = \pm 15$ equally likely, and $\mathit{MD}(E) = 6E$, find the loss of a tax and of a quota in each case.
Which is better?
\end{exercise}

\begin{exercise}[$\star\star$]\label{exo:b2:externalities-and-public-goods-theory:9}
Read the record on bees. Why do payments run from growers to beekeepers for cherries and the other way for mint? Why is a
custom used among neighbouring orchards rather than contracts?
\end{exercise}

\begin{exercise}[$\star\star\star$]\label{exo:b2:externalities-and-public-goods-theory:10}
With linear marginal benefits $a_i - G$ and marginal cost $c$, show that in the \omterm{def:b2:static-games:nash-equilibrium}{Nash equilibrium} of voluntary purchases only the
household with the largest $a_i$ buys, and that $G = \max(a_i - c, 0)$. For $n$ identical households with $a_i = a > c$, show that the gap between the
efficient and the voluntary quantity rises with $n$.
\end{exercise}

\begin{exercise}[$\star\star\star$]\label{exo:b2:externalities-and-public-goods-theory:11}
In the river of the weekend problem, a regulator believes marginal damage is $3E$ when it is $1.5E$. What quota does it set, and
what surplus is lost? Compare with the loss of no regulation.
\end{exercise}

\begin{exercise}[$\star\star\star$]\label{exo:b2:externalities-and-public-goods-theory:12}
In the model of \cref{prop:b2:externalities-and-public-goods-theory:slopes}, show that a tax reaches the efficient level for every $\theta$
when marginal damage is constant, and that a quota does when marginal damage becomes vertical at a threshold. Interpret.
\end{exercise}

\section{Problem: Brenmoor's river}\label{sec:b2:externalities-and-public-goods-theory:problem}

\begin{problem}[{Weekend problem --- a paper mill on Brenmoor's river discharges waste that harms a fishery downstream: what
level of discharge is efficient, what do a tax, permits and bargaining achieve, and which instrument is better when the
mill's costs are uncertain}]
\label{pb:b2:externalities-and-public-goods-theory:1}
The mill's marginal benefit from discharging $E$ tonnes a month is $90 - 3E$ thousand dalens a tonne; the fishery's marginal
damage is $1.5E$. Utilities are quasi-linear.

\medskip
\noindent\textbf{Part I --- A mill and a fishery.}
\begin{enumerate}
  \item How much does the mill discharge unregulated?
  \item What is the mill's total benefit from that discharge?
  \item What is the fishery's total damage?
  \item What is the net surplus?
  \item Which effect is external, and why does the mill ignore it?
\end{enumerate}

\medskip
\noindent\textbf{Part II --- Efficient discharge and bargaining.}
\begin{enumerate}[resume]
  \item Find the efficient discharge.
  \item Find the net surplus there.
  \item What surplus does the unregulated mill waste?
  \item If the mill holds the right to discharge, what does it lose and the fishery gain from a cut to the efficient level?
    What payments would both accept?
  \item If the fishery holds the right to clean water, what payments for permission to discharge would both accept?
  \item Show that the discharge agreed is the same in both cases.
  \item Why might bargaining fail on a real river?
\end{enumerate}

\medskip
\noindent\textbf{Part III --- A tax and permits.}
\begin{enumerate}[resume]
  \item What tax per tonne reaches the efficient level?
  \item What revenue does it raise?
  \item What is the mill's surplus under the tax?
  \item If the town auctions permits for the efficient total, what is their price?
  \item If it gives the permits to the mill free, what changes?
  \item A regulator who believes the damage is $3E$ sets a quota. What quota, and what surplus is lost?
\end{enumerate}

\medskip
\noindent\textbf{Part IV --- Uncertain costs: tax or quota.} The mill's marginal benefit is $90 + \theta - 3E$, with $\theta = 15$ or $-15$, equally likely; the
town must fix a tax of 30 or a quota of 20 before $\theta$ is known.
\begin{enumerate}[resume]
  \item Find the efficient discharge in each case.
  \item Find the discharge under the tax in each case.
  \item Find the loss of the tax in each case.
  \item Find the loss of the quota in each case.
  \item Which instrument has the lower expected loss? Check with the slopes rule.
  \item If the damage were $4.5E$, with the tax and quota set for the average case, which would be better?
  \item State the result: the efficient level, the \omterm{def:b2:externalities-and-public-goods-theory:corrective-tax}{corrective tax}, and the expected welfare loss of each instrument under
    cost uncertainty.
\end{enumerate}
\end{problem}
